\section{未来展望与 AI 前沿}

\subsection{RL for LLM 还需要 Breakthrough 吗？}

翁家翌的回答是\textbf{有可能}，但更紧迫的任务是先把现有的方法和计算资源榨干。当前的状态是"还没有 scale up 完全"——先 hill climb，看看现有方法的上限在哪里，再考虑是否需要新范式。

他强调了一个关键认知：\textbf{不能用当前的状态来预测接下来会发生什么}。可能有新的 RL 范式，可能有新的 post-training 范式，都有可能。每天都面对未知的挑战。

\subsection{最大的瓶颈在哪？}

翁家翌认为未来大模型的最大瓶颈仍然是 infra：

\begin{enumerate}
    \item 修 infra 的\textbf{吞吐量}——单位时间内能修多少 bug
    \item 单位时间内能\textbf{正确迭代}多少次
    \item 这两个指标决定了整个团队的生产效率，可以赋能所有其他工作
\end{enumerate}

"算法也好，环境也好——如果你把 bug 全修了，有可能算法连改都不用改就很好了。"

\subsection{如果 AI 能解决一个世界难题}

被问到希望 AI 解决什么世界难题时，翁家翌的回答出人意料：\textbf{如何预测未来}——不是预测杯子怎么掉，而是预测整个人生、世界格局、所有的一切。这个回答与他的决定论世界观一脉相承——如果世界是确定性的马尔科夫过程，理论上就应该可以被预测。

但他同时也意识到这样的工具可能是一场灾难："如果拿到一个能够预测未来的机器，那对个人而言其实是一个灾难——会导致所有价值体系的崩塌。"他甚至认为，如果有这样的 AI 模型被开发出来，"最好的选择是毁掉它，让它永远不要出来"。

\subsection{翁家翌的日常与身体管理}

翁家翌透露了他在 OpenAI 的工作状态。他认为自己的日常维护工作"并不需要那么多智商"——只要在对的方向上做对的事就好。他还分享了一个反差感很强的细节：在清华时体育课 3000 米不及格，但现在养成了每周两次跑 3000 米的习惯——"首先你要确保你的身体是健康的"。这是他"投资未来"哲学在身体管理上的延伸。

\begin{warningbox}{天才的谦逊与自我认知}
翁家翌多次强调自己并不特别——"如果把我换做任何一个人，如果他有我的 context 的话，他应该也完全可以胜任"。他认为自己很幸运在这个位置，但这个事情"换任何一个正常的人类也可以做"。这种谦逊可能是真诚的——他把成功更多归因于\textbf{正确的位置}（OpenAI + RL Infra）和\textbf{充足的 context}，而非个人天赋。
\end{warningbox}

\subsection{本章小结}

在翁家翌看来，AI 的未来既充满确定性（scale up 的路径清晰，infra 需要不断优化），也充满不确定性（新范式随时可能出现）。他选择的应对策略一如既往——在对的方向上做对的事，把 infra 打好，让快速迭代来解决一切。


